%
% beadle-identity.tex
%
% Z Specification for Beadle Identity, Address Book, and Permissions
% Formal model of the identity resolution chain, contact management,
% and per-identity rwx permission system.
%

\documentclass[a4paper,10pt,fleqn]{article}
\usepackage[margin=1in]{geometry}
\usepackage{fuzz}
\usepackage[colorlinks=true,linkcolor=blue,citecolor=blue,urlcolor=blue]{hyperref}

\begin{document}

\title{Beadle Identity and Permissions: A Z Specification}
\author{Formal Model of Identity Resolution, Address Book, and rwx Permissions}
\date{March 2026}
\hypersetup{
  pdftitle={Beadle Identity and Permissions: A Z Specification},
  pdfauthor={Punt Labs},
  pdfsubject={Z Specification},
  pdfcreator={fuzz/probcli}
}
\maketitle

\tableofcontents
\newpage

\section{Introduction}

Beadle's identity and permission systems govern three concerns:

\begin{enumerate}
\item \textbf{Identity resolution}: determining which email address beadle
  operates as, via a chain of sources (ethos, default file).
\item \textbf{Address book}: a contact store with uniqueness constraints on
  names and aliases, supporting name-based address resolution.
\item \textbf{Permissions}: an rwx bit model per (identity, contact) pair
  that gates what beadle may do with messages from each contact.
\end{enumerate}

\section{Basic Types}

\begin{zed}
[EMAIL, NAME, HANDLE, GPGKEYID]
\end{zed}

$EMAIL$ represents email addresses (e.g., \texttt{claude@punt-labs.com}).
$NAME$ represents display names and aliases (case-insensitive in the
implementation, modeled here as distinct values).
$HANDLE$ represents ethos identity handles (e.g., \texttt{claude}).
$GPGKEYID$ represents GPG key identifiers.

\section{Free Types}

\subsection{Boolean}

\begin{zed}
ZBOOL ::= ztrue | zfalse
\end{zed}

\subsection{Identity Source}

The source from which the active identity was resolved, ordered by
priority (ethos highest, defaultFile lowest).

\begin{zed}
IdentitySource ::= ethos | defaultFile
\end{zed}

\subsection{Trust Level}

Transport-level trust classification for incoming messages.

\begin{zed}
TrustLevel ::= trusted | verified | untrusted | unverified
\end{zed}

\section{Identity}

An identity is the email address beadle operates as, plus optional
metadata from the resolution source.

\begin{schema}{Identity}
email : EMAIL \\
handle : HANDLE \\
name : NAME \\
gpgKeyId : GPGKEYID \\
source : IdentitySource \\
hasHandle : ZBOOL \\
hasGpgKey : ZBOOL
\end{schema}

\subsection{Identity Resolution Chain}

The resolver checks sources in priority order. At most one identity
is active at any time.

\begin{schema}{IdentityStore}
activeIdentity : EMAIL \\
identitySource : IdentitySource \\
ethosIdentities : HANDLE \pfun EMAIL \\
defaultIdentity : EMAIL \\
hasEthos : ZBOOL \\
hasDefault : ZBOOL
\where
hasEthos = ztrue \implies activeIdentity \in \ran ethosIdentities \\
hasEthos = ztrue \implies identitySource = ethos \\
hasEthos = zfalse \land hasDefault = ztrue \implies
  activeIdentity = defaultIdentity \\
hasEthos = zfalse \land hasDefault = ztrue \implies
  identitySource = defaultFile
\end{schema}

\begin{schema}{InitIdentityStore}
IdentityStore'
\where
ethosIdentities' = \emptyset \\
hasEthos' = zfalse \\
hasDefault' = zfalse \\
hasLegacy' = zfalse
\end{schema}

\section{Permissions}

Permissions are a top-level entity: a relation between identities and
contacts, not a property of a contact. This separates access control
from contact data and simplifies permission updates.

\subsection{Permission Bits}

Each permission is a triple of read/write/execute bits.

\begin{schema}{Permission}
canRead : ZBOOL \\
canWrite : ZBOOL \\
canExecute : ZBOOL
\end{schema}

The default permission is no access (whitelist model).

\begin{axdef}
noPermission : ZBOOL \cross ZBOOL \cross ZBOOL
\where
noPermission = (zfalse, zfalse, zfalse)
\end{axdef}

\subsection{Permission Semantics}

Permission bits control what beadle may do with a contact's messages:

\begin{itemize}
\item \textbf{r} (read): beadle may read the message content
\item \textbf{w} (write): the LLM may reply and generate text (no tool use)
\item \textbf{x} (execute): beadle may execute tools as part of a reply
\end{itemize}

\subsection{Permission Store}

The permission store maps (identity email, contact name) pairs to
permission triples. It is independent of the address book.

\begin{schema}{PermissionStore}
perms : (EMAIL \cross NAME) \pfun (ZBOOL \cross ZBOOL \cross ZBOOL)
\end{schema}

\begin{schema}{InitPermissionStore}
PermissionStore'
\where
perms' = \emptyset
\end{schema}

\section{Address Book}

\subsection{Contact}

\begin{schema}{Contact}
contactName : NAME \\
contactEmail : EMAIL \\
aliases : \power NAME \\
gpgKeyId : GPGKEYID \\
hasGpgKey : ZBOOL
\where
contactName \notin aliases
\end{schema}

A contact's name is not in its alias set (they are distinct identifiers).

\subsection{Address Book State}

The address book maps names to contacts, enforcing uniqueness of all
identifiers (names and aliases) across the entire book.

\begin{schema}{AddressBook}
contacts : NAME \pinj Contact \\
allNames : \power NAME \\
identityEmail : EMAIL
\where
\forall n : \dom contacts @ \\
\quad~ (contacts~n).contactName = n \\
allNames = \dom contacts \cup
  \bigcup \{ n : \dom contacts @ (contacts~n).aliases \} \\
\forall n1, n2 : \dom contacts | n1 \neq n2 @ \\
\quad~ (\{ n1 \} \cup (contacts~n1).aliases) \cap
  (\{ n2 \} \cup (contacts~n2).aliases) = \emptyset
\end{schema}

Key invariants:
\begin{itemize}
\item Contact names are unique (partial injection).
\item The union of all names and aliases is disjoint across contacts.
\item Each contact's key matches its internal name field.
\end{itemize}

\subsection{Initialization}

\begin{schema}{InitAddressBook}
AddressBook'
\where
contacts' = \emptyset \\
allNames' = \emptyset
\end{schema}

\section{Operations}

\subsection{Add Contact}

Adding a contact succeeds only if the name and all aliases are fresh
(no conflicts with any existing contact's name or aliases).

\begin{schema}{AddContact}
\Delta AddressBook \\
newName? : NAME \\
newEmail? : EMAIL \\
newAliases? : \power NAME
\where
newName? \notin allNames \\
newAliases? \cap allNames = \emptyset \\
newName? \notin newAliases? \\
\exists c : Contact @ \\
\quad~ c.contactName = newName? \land \\
\quad~ c.contactEmail = newEmail? \land \\
\quad~ c.aliases = newAliases? \land \\
\quad~ contacts' = contacts \cup \{ newName? \mapsto c \} \\
allNames' = allNames \cup \{ newName? \} \cup newAliases? \\
identityEmail' = identityEmail
\end{schema}

\subsection{Remove Contact}

\begin{schema}{RemoveContact}
\Delta AddressBook \\
targetName? : NAME
\where
targetName? \in \dom contacts \\
contacts' = \{ targetName? \} \ndres contacts \\
allNames' = allNames \setminus
  (\{ targetName? \} \cup (contacts~targetName?).aliases) \\
identityEmail' = identityEmail
\end{schema}

\subsection{Find Contact}

Finding a contact by query matches against name, email, or any alias.
This is a query operation (no state change).

\begin{schema}{FindContact}
\Xi AddressBook \\
query? : NAME \\
results! : \power Contact
\where
results! = \{ n : \dom contacts |
  (contacts~n).contactName = query? \lor \\
\quad~ query? \in (contacts~n).aliases @ contacts~n \}
\end{schema}

\subsection{Resolve Address}

Resolve a name to an email. Succeeds only when exactly one contact
matches (no ambiguity).

\begin{schema}{ResolveAddress}
\Xi AddressBook \\
token? : NAME \\
resolved! : EMAIL
\where
\exists_1 n : \dom contacts @
  (contacts~n).contactName = token? \lor
  token? \in (contacts~n).aliases \\
\forall n : \dom contacts |
  (contacts~n).contactName = token? \lor
  token? \in (contacts~n).aliases @
  resolved! = (contacts~n).contactEmail
\end{schema}

\subsection{Check Permission}

Look up the effective permission for a contact given the active identity.
Returns the explicit permission if set, otherwise no-permission (whitelist).

\begin{schema}{CheckPermission}
\Xi PermissionStore \\
identityEmail? : EMAIL \\
targetName? : NAME \\
permRead! : ZBOOL \\
permWrite! : ZBOOL \\
permExecute! : ZBOOL
\where
((identityEmail?, targetName?) \in \dom perms \implies \\
\quad~ (permRead!, permWrite!, permExecute!) =
  perms~(identityEmail?, targetName?)) \\
((identityEmail?, targetName?) \notin \dom perms \implies \\
\quad~ (permRead!, permWrite!, permExecute!) = noPermission)
\end{schema}

\subsection{Set Permission}

Grant or change permission for a contact.

\begin{schema}{SetPermission}
\Delta PermissionStore \\
forIdentity? : EMAIL \\
targetName? : NAME \\
newRead? : ZBOOL \\
newWrite? : ZBOOL \\
newExecute? : ZBOOL
\where
perms' = perms \oplus
  \{ (forIdentity?, targetName?) \mapsto
    (newRead?, newWrite?, newExecute?) \}
\end{schema}

\section{Unified Initialization}

\begin{schema}{Init}
InitAddressBook \\
InitIdentityStore \\
InitPermissionStore
\end{schema}

\section{System Invariants}

Key properties that hold across all reachable states:

\begin{enumerate}
\item \textbf{Name uniqueness}: no two contacts share a name or alias.
  Enforced by the partial injection $contacts$ and the disjointness
  constraint on $allNames$.
\item \textbf{Whitelist default}: absent permission entries default to
  $---$ (no access). There is no implicit grant.
\item \textbf{Identity resolution is deterministic}: exactly one source
  provides the active identity, selected by priority chain.
\item \textbf{Fail-closed resolution}: if ethos has an active handle but
  the identity file is corrupt, resolution fails rather than falling
  back to a lower-priority source.
\end{enumerate}

\end{document}
